\documentclass[tikz,border=12pt]{standalone}
\usepackage[T1]{fontenc}
\usepackage{mathpazo}
\usepackage{amsmath,amssymb}
\usetikzlibrary{arrows.meta, positioning, calc, decorations.pathreplacing}

\definecolor{stblue}{HTML}{7B97AD}
\definecolor{rwgold}{HTML}{B8995E}
\definecolor{sage}{HTML}{7A9A7E}
\definecolor{dkbrown}{HTML}{3D3530}
\definecolor{wmgray}{HTML}{7B7067}
\definecolor{cream}{HTML}{F9F6F0}
\definecolor{border}{HTML}{D4C9B8}
\definecolor{terracotta}{HTML}{C47A6A}
\definecolor{lavender}{HTML}{907EA0}

\begin{document}
\begin{tikzpicture}[>={Stealth[length=4mm, width=3mm]}, line width=1.5pt]

% Background
\fill[cream, rounded corners=14pt] (-1.5,-2.5) rectangle (13,8.5);
\draw[border, line width=1.5pt, rounded corners=14pt] (-1.5,-2.5) rectangle (13,8.5);

% Title
\node[font=\LARGE\bfseries, text=dkbrown] at (5.75,7.8)
  {Descent Direction};

% === Contour plot of f (elliptical contours) ===

% Draw several contour ellipses (tilted slightly)
\foreach \r/\op in {1.2/0.06, 2.0/0.06, 2.8/0.06, 3.6/0.06, 4.4/0.06} {
  \draw[stblue!\op!stblue, line width=0.8pt, stblue!25]
    (7.5,3.0) ellipse ({\r*1.4} and {\r*0.7});
}

% Fill innermost contour
\fill[stblue!8] (7.5,3.0) ellipse (1.68 and 0.84);

% Redraw contours on top
\foreach \r in {1.2, 2.0, 2.8, 3.6, 4.4} {
  \draw[stblue!30, line width=0.8pt]
    (7.5,3.0) ellipse ({\r*1.4} and {\r*0.7});
}

% Minimum point x*
\fill[stblue] (7.5,3.0) circle (4pt);
\node[font=\large\bfseries, text=stblue, below] at (7.5,2.7) {$x^*$};

% Current point x
\fill[terracotta] (3.5,4.2) circle (5pt);
\node[font=\Large\bfseries, text=terracotta, above left] at (3.5,4.4) {$x$};

% Gradient direction (pointing away from minimum)
\draw[->, terracotta!60, line width=2pt] (3.5,4.2) -- (1.8,4.8);
\node[font=\large\bfseries, text=terracotta!70, above] at (1.8,5.0) {$\nabla f(x)$};

% Negative gradient (descent direction)
\draw[->, sage, line width=2.5pt] (3.5,4.2) -- (5.2,3.6);
\node[font=\large\bfseries, text=sage!80!black, below right] at (5.0,3.4)
  {$-\nabla f(x)$};

% General descent direction d (at angle)
\draw[->, rwgold, line width=2.5pt] (3.5,4.2) -- (5.8,4.6);
\node[font=\large\bfseries, text=rwgold, above] at (5.8,4.8) {$d$};

% Dashed arc showing the cone of descent directions
% The gradient points roughly at angle 160 deg from x, so descent directions
% are those with angle < 90 from -gradient direction
\draw[wmgray, line width=1pt, dashed]
  (3.5,4.2) ++ (70:2.2) arc (70:-50:2.2);
\node[font=\large, text=wmgray, align=center] at (6.2,6.4)
  {descent directions:\\$\nabla f(x)^\top d < 0$};

% Separating line (perpendicular to gradient through x)
% Gradient direction is roughly (1.8-3.5, 4.8-4.2) = (-1.7, 0.6)
% Perpendicular is (0.6, 1.7)
\draw[wmgray, line width=0.8pt, dotted]
  ($(3.5,4.2) + (0.6,1.7)*1.2$) -- ($(3.5,4.2) - (0.6,1.7)*1.2$);

% Annotation: condition
\node[font=\large, text=dkbrown, fill=cream, inner sep=4pt,
      rounded corners=3pt, align=left] at (5.75,-1.5)
  {$d$ is a descent direction at $x$ if
   $\nabla f(x)^\top d < 0$
   (acute angle with $-\nabla f(x)$)};

% Small angle arc between d and -grad f
\draw[rwgold, line width=1pt, ->] ($(3.5,4.2) + (-20:1.2)$) arc (-20:12:1.2);
\node[font=\small, text=rwgold] at (4.9,3.9) {$\theta$};

\end{tikzpicture}
\end{document}
